\Article{8}
Wir wollen nun unsere allgemeine Lösung noch an einigen Beispielen erläutern. So lässt sich am besten zeigen, wie man die Lösung anwendet und welcher Art einige weitere dabei zu berücksichtigende Umstände sind.

Zunächst seien beide Flächen Ebenen. Dann können wir
\[
\begin{aligned}
x&=t,&y&=u,&z&=0\\
X&=T,&Y&=U,&Z&=0
\end{aligned}
\]
\noindent setzen. Die Differentialgleichung
\[
\omega=\mathrm{d}t^2+\mathrm{d}u^2=0
\]
\noindent liefert hier die beiden Integrale
\[
t+iu=\text{Const.},\quad t-iu=\text{Const.}
\]
\noindent Ebenso lauten die beiden Integrale der Gleichung \(\Omega=\mathrm{d}T^2+\mathrm{d}U^2=0\)
\[
T+iU=\text{Const.},\quad T-iU=\text{Const.}
\]
\noindent Die beiden allgemeinen Lösungen der Aufgabe sind demnach
\[
\begin{aligned}
\mathrm{I.}\quad&T+iU=f(t+iu),\quad T-iU=f'(t-iu)\\
\mathrm{II.}\quad&T+iU=f(t-iu),\quad T-iU=f'(t+iu).
\end{aligned}
\]
\par\indent\SourceLeaf{205}%
Dieses Ergebnis lässt sich auch so ausdrücken. Wenn das Zeichen \(f\) für eine beliebige Funktion steht, nimmt man den Realteil von \(f(x+iy)\) für \(X\) und den Imaginärteil unter Weglassung des Faktors \(i\) entweder für \(Y\) oder für \(-Y\).

Verwendet man die Zeichen \(\varphi,\varphi'\) in der Bedeutung von Artikel 7 und setzt
\[
\varphi(x+iy)=\xi+i\eta,\quad\varphi'(x-iy)=\xi-i\eta
\]
\noindent wobei \(\xi\) und \(\eta\) offenbar reelle Funktionen von \(x\) und \(y\) sind, so erhält man in der \textit{ersten} Lösung
\[
\begin{aligned}
\mathrm{d}X+i\,\mathrm{d}Y&=(\xi+i\eta)(\mathrm{d}x+i\,\mathrm{d}y)\\
\mathrm{d}X-i\,\mathrm{d}Y&=(\xi-i\eta)(\mathrm{d}x-i\,\mathrm{d}y)
\end{aligned}
\]
\noindent und folglich
\[
\begin{aligned}
\mathrm{d}X&=\xi\,\mathrm{d}x-\eta\,\mathrm{d}y\\
\mathrm{d}Y&=\eta\,\mathrm{d}x+\xi\,\mathrm{d}y
\end{aligned}
\]
\noindent Setzt man nun
\[
\begin{aligned}
\xi&=\sigma\mathord{.}\cos\gamma,&\eta&=\sigma\mathord{.}\sin\gamma\\
\mathrm{d}x&=\mathrm{d}s\mathord{.}\cos g,&\mathrm{d}y&=\mathrm{d}s\mathord{.}\sin g\\
\mathrm{d}X&=\mathrm{d}S\mathord{.}\cos G,&\mathrm{d}Y&=\mathrm{d}S\mathord{.}\sin G
\end{aligned}
\]
\noindent sodass \(\mathrm{d}s\) ein Linienelement in der ersten Ebene, \(g\) seine Neigung gegen die Abszissenachse, \(\mathrm{d}S\) das entsprechende Linienelement in der zweiten Ebene und \(G\) seine Neigung gegen die Abszissenachse bezeichnet, so liefern die obigen Gleichungen
\[
\begin{aligned}
\mathrm{d}S\mathord{.}\cos G&=\sigma\mathord{.}\mathrm{d}s\mathord{.}\cos(g+\gamma)\\
\mathrm{d}S\mathord{.}\sin G&=\sigma\mathord{.}\mathrm{d}s\mathord{.}\sin(g+\gamma)
\end{aligned}
\]
\noindent und folglich, wenn man. was zulässig ist, \(\sigma\) als positiv betrachtet,
\[
\mathrm{d}S=\sigma\mathord{.}\mathrm{d}s,\quad G=g+\gamma
\]
\noindent Man sieht also (in Übereinstimmung mit Artikel 7), dass \(\sigma\) das Vergrößerungsverhältnis des Elements \(\mathrm{d}s\) in seiner Abbildung \(\mathrm{d}S\) angibt und. wie es sein soll, von \(g\) unabhängig ist. Ebenso zeigt die Unabhängigkeit des Winkels \(\gamma\) von \(g\), dass alle von einem Punkt ausgehenden Linienelemente in der ersten Ebene
\SourceLeaf{206}
durch Elemente in der zweiten Ebene abgebildet werden, die miteinander dieselben Winkel bilden wie jene, und zwar, wie wir hinzufügen können, \textit{im gleichen Sinn}.

Wählt man für \(f\) eine lineare Funktion, sodass \(f\upsilon=A+B\upsilon\) gilt, wobei die konstanten Koeffizienten die Form
\[
A=a+bi,\quad B=c+ei
\]
\noindent haben, so erhält man
\[
\varphi\upsilon=B=c+ei
\]
\noindent und somit
\[
\sigma=\BareRoot{(cc+ee)},\quad\gamma=\mathrm{Arc.}\,\mathrm{tang}\frac{e}{c}
\]
\noindent Das Vergrößerungsverhältnis ist folglich in allen Punkten konstant, und die Abbildung ist dem Abgebildeten in jeder Hinsicht ähnlich.

Für jede andere Funktion \(f\) ist das Vergrößerungsverhältnis, wie sich leicht zeigen lässt, veränderlich. Die Ähnlichkeit kann daher nur in den kleinsten Teilen bestehen.

Sind für eine bestimmte Anzahl gegebener Punkte der ersten Ebene die entsprechenden Positionen in der Abbildung vorgegeben, so lässt sich mit der üblichen Interpolationsmethode leicht die einfachste algebraische Funktion \(f\) finden, die diese Bedingung erfüllt. Bezeichnet man die Werte von \(x+iy\) an den gegebenen Punkten mit \(a,b,c\) usw. und die entsprechenden Werte von \(X+iY\) mit \(A,B,C\) usw., so muss man
\[
\begin{aligned}
f\upsilon={}&\frac{(\upsilon-b)(\upsilon-c).\ .\ .}{(a-b)(a-c).\ .\ .}\mathord{.}A
+\frac{(\upsilon-a)(\upsilon-c).\ .\ .}{(b-a)(b-c).\ .\ .}\mathord{.}B\\
&+\frac{(\upsilon-a)(\upsilon-b).\ .\ .}{(c-a)(c-b).\ .\ .}\mathord{.}C+\text{etc.}
\end{aligned}
\]
\noindent setzen. Das ist eine algebraische Funktion von \(\upsilon\), deren Grad um eins kleiner ist als die Anzahl der vorgegebenen Punkte. Bei zwei Punkten wird die Funktion linear, sodass vollständige Ähnlichkeit besteht.

Dieses Verfahren ist in der Geodäsie nützlich. Kennt man die richtige Lage einer Anzahl von Punkten, so lässt sich damit eine Karte verbessern, die auf mittelmäßigen Messungen beruht und in den feinen Einzelheiten gut, im Ganzen aber etwas verzerrt ist. Es versteht sich jedoch, dass man bei einer solchen Umformung nicht weit über das Gebiet hinausgehen darf, das diese Punkte umschließen.

Führt man die \textit{zweite} Lösung auf dieselbe Weise durch, so zeigt sich als einziger Unterschied, dass die Ähnlichkeit seitenverkehrt ist. Die Elemente in der Abbildung bilden miteinander zwar ebenso große Winkel wie im Abgebildeten, aber im entgegengesetzten Sinn, sodass
\SourceLeaf{207}
dort rechts liegt, was hier links liegt. Dieser Unterschied ist jedoch unwesentlich und verschwindet, wenn man in einer der Ebenen die zuvor als obere betrachtete Seite zur unteren macht. Diese letzte Bemerkung gilt immer, wenn eine der beiden Flächen eine Ebene ist. Deshalb können wir uns in den folgenden Beispielen dieser Art auf die erste Lösung beschränken.

